\section{可检验性、估计误差与时间外推}
\label{sec:verification}

主判据是有限的，但“有限”只保证问题有明确定义并能穷举终止，不保证计算高效、统计稳定或长时间有效。本节分别说明这三条边界。

\subsection{固定有限系统上怎样检验}

对固定有限 $B$、有限 rational kernel 族和给定任务，可以枚举全部任务保真分区，逐个计算 $D_{\bar\cK}(\pi)$，从而精确求出复杂度--直径前沿与 $\Gamma$。固定分区的最优宏观核也可写成有限维线性规划。分区数按 Bell 数增长，所以这种程序只保证可终止，不保证实际规模上的效率。

若核含一般可计算实数，只能把结果逼近到给定精度；若 $\bar\cK$ 是无限族，还需要紧致连续参数化、可控有限覆盖、解析上界或相应 oracle，才能处理对 $\bar K$ 的上确界。对有限多个规模的数值结果，也不能单独证明 $N\to\infty$ 的渐近命题。结构证书的价值正在于：它们用对称性、充分统计量或模型专属估计绕开完整分区枚举，但仍须给出主判据所需的两个上界。

\subsection{固定操作性空间后的估计稳健性}

设真实与估计诱导核按同一参数 $\theta$ 配对，并在已经固定的 $B_N$ 上满足
\[
\sup_\theta d_\infty(\bar K_{N,\theta},\widehat{\bar K}_{N,\theta})
\le\delta_N.
\]

\begin{proposition}[证书扰动界]
\label{prop:estimation-robustness}
对每个候选分区 $\pi$，
\[
\bigl|D_{\bar\cK_N}(\pi)-D_{\widehat{\bar\cK}_N}(\pi)\bigr|
\le2\delta_N,
\]
并且
\[
|\Gamma_N-\widehat\Gamma_N|\le2\delta_N.
\]
\end{proposition}

\begin{proof}
确定性推前不增加 TV 距离，所以每个投影后的行分布改变至多 $\delta_N$；一对行的距离因此改变至多 $2\delta_N$。依次取最大值、对核族取上确界，仍保持该界；最后对任务保真分区取最小值，得到 $\Gamma_N$ 的同一扰动界。
\end{proof}

这里“$B_N$ 已经固定”是实质条件。操作性微观分区由精确等式定义，任意小的 kernel 扰动都可能使它细化；上面的命题不保证从有限样本稳定恢复 $B_N$ 本身。若微观基线未知，实证工作需要另行声明近似等价阈值、统计间隔或测量分辨率。

\subsection{一步闭合怎样延伸到时间轴}

设某个候选 $Q$ 满足一步闭合缺陷
\[
d_\infty(KQ,QL)\le\varepsilon.
\]
需要区分时刻 $T$ 的\emph{终端边缘分布}与从 $0$ 到 $T$ 的\emph{完整宏观路径分布}。

对终端边缘，递推与 TV 非扩张给出
\begin{equation}
d_\infty(K^TQ,QL^T)
\le
\varepsilon\sum_{j=0}^{T-1}\alpha(L^j)
\le T\varepsilon.
\label{eq:body-endpoint-bound}
\end{equation}
因此，若 $\varepsilon_N\to0$，任意固定 $T$ 的终端边缘都收敛；当
$Q_N=Q_{q_N}$ 来自确定性粗粒化时，下述耦合还给出固定长度的投影路径律收敛。
若时间窗 $T_N$ 增长，终端边缘的直接充分条件是
\[
\varepsilon_N\sum_{j<T_N}\alpha(L_N^j)\longrightarrow0.
\]
特别地，若 $\alpha(L_N)\le\rho<1$ 一致成立，则终端误差至多 $\varepsilon_N/(1-\rho)$，可对任意 $T_N$ 一致控制。

完整路径的要求更强。若 $Q=Q_q$ 是确定性粗粒化，逐步耦合给出的通用界是
\begin{equation}
d_{\TV}(\text{两种长度 }T\text{ 的宏观路径律})
\le1-(1-\varepsilon)^T
\le T\varepsilon.
\label{eq:body-path-bound}
\end{equation}
所以 $T_N\varepsilon_N\to0$ 是增长路径窗的直接充分条件。宏观核的收缩可以淡化当前状态的旧误差，却不能抹去路径中已经记录的早期偏差；因此不能用终端边缘的统一收缩界替代完整路径条件。$\Gamma_N\to0$ 提供的是某列候选的一步误差趋零，它本身不指定增长时间窗、收缩率或路径耦合。完整推导和近似证书的空间复合见附录~\ref{app:composition-time}。

\subsection{解释主定理时必须保留的边界}

\paperclaim{量词边界。}
正文允许每个 $\bar K$ 选择自己的最优 $L$。若现实问题要求不同环境共享同一宏观动力学，必须研究更强的 $\min_L\sup_{\bar K}$；二者一般不能交换。附录~\ref{app:examples-counterexamples}给出有限反例。

\paperclaim{复杂度边界。}
状态类别少不等于描述短、参数少、求值快、样本复杂度低、语义清晰或几何维数低。若应用需要其中任何一种复杂度，必须另行定义并证明相应判据。

\paperclaim{任务与因果边界。}
任务、探针和允许动力学先于定理给定；改变其中任何一项都可能改变最粗商和前沿。本文的自治只表示声明动力学下的概率闭合，不自动表示干预意义上的因果自治，也不把任务成功等同于机制识别。

\paperclaim{分布边界。}
最坏行误差与分布加权平均误差回答不同问题。后者可以忽略零质量状态并容许低概率坏状态；没有覆盖条件时，不能把加权证书直接当作 $d_\infty$ 证书。
